\section{Arm Cortex-M4 \& Timer}

\subsection{Watchdog Window-Watchdog-Timer (WWDT)}
\begin{itemize}
    \item Zur Überwachung der Code-Ausführung. Erzeugt im Fehlerfall einen Reset.
    \item Dieser Timer muss jeweils vor Ablauf der Timeout-Zeit zurückgesetzt werden.
    \item Ansonsten wird eine definierte Aktion ausgeführt (z.B. Reset).
\end{itemize}

\subsection{Spezialtimer STM32G4}
\begin{itemize}
    \item HRTIM (high resolution)
    \item LPTIM (low power)
    \item RTC
    \item WWDG (eigenen Code überwachen)
\end{itemize}

\subsection{M4 vs M0}
\begin{itemize}
    \item \textbf{M4:} schnellere Taktung als M0
    \item \textbf{M0:} Low Power / Low Cost (alte Technologie)
    \item \textbf{M4:} High Performance / neuere Technologie
\end{itemize}

\subsubsection{Cortex-M4 Prozessor}
\begin{itemize}
    \item 3-stufige Pipeline für Thumb-2 Instruktionssatz
    \item Erweiterte Busstruktur mit Harvard-Architektur (I-Code, D-Code, System-Bus \& PPB)
    \item NVIC, SYSTICK \& leistungsfähige debug Units
    \item Memory Protection Unit (MPU)
    \item Floating Point Unit (FPU)
\end{itemize}

\subsubsection{Vorteile Cortex-M4 gegenüber Cortex-M0}
\begin{itemize}
    \item Floating Point Unit, erweiterte Debug-Unit
    \item Harvard-Architektur
    \item Bit-Banding, SIMD-Befehle, erweiterter Befehlssatz, DSP-Funktionen.
    \item Schnellere Ausführung von Code, unterstützt Thumb-2.
    \item ALU und Barrel Shifter und Multiply \& Accumulate Unit (MAC)
\end{itemize}

\subsection{Wichtigkeit Standardisierung von ARM-Cortex}
\begin{itemize}
    \item Verwendung des gleichen Cores durch verschiedene Chip-Hersteller.
    \item Es gibt eine Abwärtskompatibilität innerhalb der Familien.
\end{itemize}




